% Options for packages loaded elsewhere
% Options for packages loaded elsewhere
\PassOptionsToPackage{unicode}{hyperref}
\PassOptionsToPackage{hyphens}{url}
\PassOptionsToPackage{dvipsnames,svgnames,x11names}{xcolor}
%
\documentclass[
  12pt,
  letterpaper,
  DIV=11,
  numbers=noendperiod]{scrartcl}
\usepackage{xcolor}
\usepackage{amsmath,amssymb}
\setcounter{secnumdepth}{5}
\usepackage{iftex}
\ifPDFTeX
  \usepackage[T1]{fontenc}
  \usepackage[utf8]{inputenc}
  \usepackage{textcomp} % provide euro and other symbols
\else % if luatex or xetex
  \usepackage{unicode-math} % this also loads fontspec
  \defaultfontfeatures{Scale=MatchLowercase}
  \defaultfontfeatures[\rmfamily]{Ligatures=TeX,Scale=1}
\fi
\usepackage{lmodern}
\ifPDFTeX\else
  % xetex/luatex font selection
  \setmainfont[]{Times New Roman}
\fi
% Use upquote if available, for straight quotes in verbatim environments
\IfFileExists{upquote.sty}{\usepackage{upquote}}{}
\IfFileExists{microtype.sty}{% use microtype if available
  \usepackage[]{microtype}
  \UseMicrotypeSet[protrusion]{basicmath} % disable protrusion for tt fonts
}{}
\makeatletter
\@ifundefined{KOMAClassName}{% if non-KOMA class
  \IfFileExists{parskip.sty}{%
    \usepackage{parskip}
  }{% else
    \setlength{\parindent}{0pt}
    \setlength{\parskip}{6pt plus 2pt minus 1pt}}
}{% if KOMA class
  \KOMAoptions{parskip=half}}
\makeatother
% Make \paragraph and \subparagraph free-standing
\makeatletter
\ifx\paragraph\undefined\else
  \let\oldparagraph\paragraph
  \renewcommand{\paragraph}{
    \@ifstar
      \xxxParagraphStar
      \xxxParagraphNoStar
  }
  \newcommand{\xxxParagraphStar}[1]{\oldparagraph*{#1}\mbox{}}
  \newcommand{\xxxParagraphNoStar}[1]{\oldparagraph{#1}\mbox{}}
\fi
\ifx\subparagraph\undefined\else
  \let\oldsubparagraph\subparagraph
  \renewcommand{\subparagraph}{
    \@ifstar
      \xxxSubParagraphStar
      \xxxSubParagraphNoStar
  }
  \newcommand{\xxxSubParagraphStar}[1]{\oldsubparagraph*{#1}\mbox{}}
  \newcommand{\xxxSubParagraphNoStar}[1]{\oldsubparagraph{#1}\mbox{}}
\fi
\makeatother


\usepackage{longtable,booktabs,array}
\newcounter{none} % for unnumbered tables
\usepackage{calc} % for calculating minipage widths
% Correct order of tables after \paragraph or \subparagraph
\usepackage{etoolbox}
\makeatletter
\patchcmd\longtable{\par}{\if@noskipsec\mbox{}\fi\par}{}{}
\makeatother
% Allow footnotes in longtable head/foot
\IfFileExists{footnotehyper.sty}{\usepackage{footnotehyper}}{\usepackage{footnote}}
\makesavenoteenv{longtable}
\usepackage{graphicx}
\makeatletter
\newsavebox\pandoc@box
\newcommand*\pandocbounded[1]{% scales image to fit in text height/width
  \sbox\pandoc@box{#1}%
  \Gscale@div\@tempa{\textheight}{\dimexpr\ht\pandoc@box+\dp\pandoc@box\relax}%
  \Gscale@div\@tempb{\linewidth}{\wd\pandoc@box}%
  \ifdim\@tempb\p@<\@tempa\p@\let\@tempa\@tempb\fi% select the smaller of both
  \ifdim\@tempa\p@<\p@\scalebox{\@tempa}{\usebox\pandoc@box}%
  \else\usebox{\pandoc@box}%
  \fi%
}
% Set default figure placement to htbp
\def\fps@figure{htbp}
\makeatother


% definitions for citeproc citations
\NewDocumentCommand\citeproctext{}{}
\NewDocumentCommand\citeproc{mm}{%
  \begingroup\def\citeproctext{#2}\cite{#1}\endgroup}
\makeatletter
 % allow citations to break across lines
 \let\@cite@ofmt\@firstofone
 % avoid brackets around text for \cite:
 \def\@biblabel#1{}
 \def\@cite#1#2{{#1\if@tempswa , #2\fi}}
\makeatother
\newlength{\cslhangindent}
\setlength{\cslhangindent}{1.5em}
\newlength{\csllabelwidth}
\setlength{\csllabelwidth}{3em}
\newenvironment{CSLReferences}[2] % #1 hanging-indent, #2 entry-spacing
 {\begin{list}{}{%
  \setlength{\itemindent}{0pt}
  \setlength{\leftmargin}{0pt}
  \setlength{\parsep}{0pt}
  % turn on hanging indent if param 1 is 1
  \ifodd #1
   \setlength{\leftmargin}{\cslhangindent}
   \setlength{\itemindent}{-1\cslhangindent}
  \fi
  % set entry spacing
  \setlength{\itemsep}{#2\baselineskip}}}
 {\end{list}}
\usepackage{calc}
\newcommand{\CSLBlock}[1]{\hfill\break\parbox[t]{\linewidth}{\strut\ignorespaces#1\strut}}
\newcommand{\CSLLeftMargin}[1]{\parbox[t]{\csllabelwidth}{\strut#1\strut}}
\newcommand{\CSLRightInline}[1]{\parbox[t]{\linewidth - \csllabelwidth}{\strut#1\strut}}
\newcommand{\CSLIndent}[1]{\hspace{\cslhangindent}#1}



\setlength{\emergencystretch}{3em} % prevent overfull lines

\providecommand{\tightlist}{%
  \setlength{\itemsep}{0pt}\setlength{\parskip}{0pt}}



 


\usepackage{tcolorbox}
\usepackage{amssymb}
\usepackage{yfonts}
\usepackage{bm}


\newtcolorbox{greybox}{
  colback=white,
  colframe=blue,
  coltext=black,
  boxsep=5pt,
  arc=4pt}
  
\newcommand{\sectionbreak}{\clearpage}

 
\newcommand{\ds}[4]{\sum_{{#1}=1}^{#3}\sum_{{#2}=1}^{#4}}
\newcommand{\us}[3]{\mathop{\sum\sum}_{1\leq{#2}<{#1}\leq{#3}}}

\newcommand{\ol}[1]{\overline{#1}}
\newcommand{\ul}[1]{\underline{#1}}

\newcommand{\amin}[1]{\mathop{\text{argmin}}_{#1}}
\newcommand{\amax}[1]{\mathop{\text{argmax}}_{#1}}

\newcommand{\ci}{\perp\!\!\!\perp}

\newcommand{\mc}[1]{\mathcal{#1}}
\newcommand{\mb}[1]{\mathbb{#1}}
\newcommand{\mf}[1]{\mathfrak{#1}}

\newcommand{\eps}{\epsilon}
\newcommand{\lbd}{\lambda}
\newcommand{\alp}{\alpha}
\newcommand{\df}{=:}
\newcommand{\am}[1]{\mathop{\text{argmin}}_{#1}}
\newcommand{\ls}[2]{\mathop{\sum\sum}_{#1}^{#2}}
\newcommand{\ijs}{\mathop{\sum\sum}_{1\leq i<j\leq n}}
\newcommand{\jis}{\mathop{\sum\sum}_{1\leq j<i\leq n}}
\newcommand{\sij}{\sum_{i=1}^n\sum_{j=1}^n}
	
\KOMAoption{captions}{tableheading}
\makeatletter
\@ifpackageloaded{caption}{}{\usepackage{caption}}
\AtBeginDocument{%
\ifdefined\contentsname
  \renewcommand*\contentsname{Table of contents}
\else
  \newcommand\contentsname{Table of contents}
\fi
\ifdefined\listfigurename
  \renewcommand*\listfigurename{List of Figures}
\else
  \newcommand\listfigurename{List of Figures}
\fi
\ifdefined\listtablename
  \renewcommand*\listtablename{List of Tables}
\else
  \newcommand\listtablename{List of Tables}
\fi
\ifdefined\figurename
  \renewcommand*\figurename{Figure}
\else
  \newcommand\figurename{Figure}
\fi
\ifdefined\tablename
  \renewcommand*\tablename{Table}
\else
  \newcommand\tablename{Table}
\fi
}
\@ifpackageloaded{float}{}{\usepackage{float}}
\floatstyle{ruled}
\@ifundefined{c@chapter}{\newfloat{codelisting}{h}{lop}}{\newfloat{codelisting}{h}{lop}[chapter]}
\floatname{codelisting}{Listing}
\newcommand*\listoflistings{\listof{codelisting}{List of Listings}}
\usepackage{amsthm}
\theoremstyle{plain}
\newtheorem{theorem}{Theorem}[section]
\theoremstyle{remark}
\AtBeginDocument{\renewcommand*{\proofname}{Proof}}
\newtheorem*{remark}{Remark}
\newtheorem*{solution}{Solution}
\newtheorem{refremark}{Remark}[section]
\newtheorem{refsolution}{Solution}[section]
\makeatother
\makeatletter
\makeatother
\makeatletter
\@ifpackageloaded{caption}{}{\usepackage{caption}}
\@ifpackageloaded{subcaption}{}{\usepackage{subcaption}}
\makeatother
\usepackage{bookmark}
\IfFileExists{xurl.sty}{\usepackage{xurl}}{} % add URL line breaks if available
\urlstyle{same}
\hypersetup{
  pdftitle={Matrix Approximation with Kronecker and Precision Weighting},
  pdfauthor={Jan de Leeuw; Jan Graffelman},
  colorlinks=true,
  linkcolor={blue},
  filecolor={Maroon},
  citecolor={Blue},
  urlcolor={Blue},
  pdfcreator={LaTeX via pandoc}}


\title{Matrix Approximation with Kronecker and Precision Weighting}
\author{Jan de Leeuw \and Jan Graffelman}
\date{April 19, 2026}
\begin{document}
\maketitle
\begin{abstract}
TBD
\end{abstract}

\renewcommand*\contentsname{Table of contents}
{
\setcounter{tocdepth}{3}
\tableofcontents
}

\sectionbreak

\textbf{Note:} This is a working manuscript which will be
expanded/updated frequently. All suggestions for improvement are
welcome. All Rmd, tex, html, pdf, R, and C files are in the public
domain. Attribution will be appreciated, but is not required. The files
can be found at \url{https://github.com/deleeuw/triple}

\sectionbreak

\section{Introduction}\label{introduction}

The problem addressed in this note is to approximate a given
\(n\times m\) matrix \(Y\) by a \(n\times m\) matrix \(Z\) which must be
in a given subset \(\mathbb{Z}\) of \(\mathbb{R}^{n\times m}\). For
example, \(\mathbb{Z}\) can be the \(n\times m\) matrices of rank less
than or equal to \(p\), or a set of special matrices, such as Hankel or
Cauchy matrices.

We measure loss using \begin{equation}
\sigma(Z):=\text{tr}\ R(W\circ(Y-Z))C(W\circ(Y-Z))'.\label{eq-loss}
\end{equation} In Definition \eqref{eq-loss} \(R\) and \(C\) are given
positive sem-definite matrices of orders \(n\) and \(m\), which define
the \emph{Kronecker weights}. The Kronecker weights can be, for example,
inverse covariance matrices for row and column entries, as in the matrix
variate normal distribution (Gupta and Nagar
(\citeproc{ref-gupta_nagar_00}{2000})). The \(n\times m\) matrix \(W\)
of \emph{precision weights} has positive entries, which can be, for
example, inverse standard errors of the corresponding entries of \(Y\).
The symbol \(\circ\) is used for Hadamard (elementwise) multiplication.

\sectionbreak

\section{Algorithm}\label{algorithm}

Our algorithm uses majorization (or MM). See Lange
(\citeproc{ref-lange_16}{2016}) for comprehensive discussion of MM.
Majorization is used to reduce the weighted least squares problem to an
unweighted least squares problem, which is generally much easier to
solve. Similar applications of MM are in Kiers
(\citeproc{ref-kiers_97}{1997}) and in Groenen et al.
(\citeproc{ref-groenen_giaquinto_kiers_03}{2003}).

We employ two nested majorizations. In the first majorization we get rid
of the Kronecker weighting and reduce each MM iteration to a weighted
least squares problem that is separable. The second majorization is used
to make a further reduction of this precision weighted least squares
problem to a sequence of unweighted least squares problems.

\begin{theorem}[]\protect\hypertarget{thm-major}{}\label{thm-major}

Suppose \(X\) is \(n\times m\) and \(R\) and \(C\) are positive
semi-definite of orders \(n\) and \(m\). Define \[
\sigma(X):=\text{tr}\ RXCX'.
\] Suppose \(\tilde X\) is any other \(n\times m\) matrix. Choose a
\(\lambda\) larger than or equal to the largest eigenvalue of the
Kronecker product \(R\otimes C\). Define
\(\tilde G:=\lambda^{-1}R\tilde XC\) and
\(\tilde V:=\tilde X-\tilde G\). Then \[
\sigma(X)\leq\sigma(\tilde X)+\lambda\|X-\tilde V\|^2-\lambda^{-1}\|\tilde G\|^2
\]

\end{theorem}

\begin{proof}
Write \(X=\tilde X+(X-\tilde X)\). Then \begin{equation}
\sigma(X)=\sigma(\tilde X)+2\text{tr}\ R(X-\tilde X)C\tilde X'+\text{tr}\ R(X-\tilde X)C(X-\tilde X)'.
\end{equation} Now, using \(x:=\text{vec}(X)\) and
\(\tilde x:=\text{vec}(\tilde X)\), \begin{multline}
\text{tr}\ R(X-\tilde X)C(X-\tilde X)'=(x-\tilde x)'C\otimes R(x-\tilde x)\leq\\\lambda(x-\tilde x)'C\otimes R(x-\tilde x)=\lambda\text{tr}\ (X-\tilde X)(X-\tilde X)',
\end{multline} and thus \begin{multline}
\sigma(X)\leq\sigma(\tilde X)+2\text{tr}\ (X-\tilde X)'\tilde G+\lambda\text{tr}\ (X-\tilde X)'(X-\tilde X)=\\\sigma(\tilde X)+\lambda\{2\text{tr}\ (X-\tilde X)'\lambda^{-1}\tilde G+\text{tr}\ (X-\tilde X)'(X-\tilde X)\}=\\
\sigma(\tilde X)+\lambda\|X-\tilde V\|^2-\lambda^{-1}\|\tilde G\|^2
\end{multline}
\end{proof}

Now suppose \(X=W\circ(Y-Z)\) and \(\tilde X=W\circ(Y-\tilde Z)\). Then
\begin{equation}
\tilde v_{ij}=w_{ij}(y_{ij}-\tilde z_{ij})-\tilde g_{ij},
\end{equation} and thus \begin{equation}
x_{ij}-\tilde v_{ij}=-w_{ij}(z_{ij}-\tilde z_{ij})+\tilde g_{ij}=
-w_{ij}(z_{ij}-(\tilde z_{ij}-w_{ij}^{-1}\tilde g_{ij})),
\end{equation} and in the majorization step we must minimize
\begin{equation}
\sum_{i=1}^n\sum_{j=1}^mw_{ij}^2(z_{ij}-(\tilde z_{ij}-w_{ij}^{-1}\tilde g_{ij}))^2
\end{equation} over \(z\). This can be done with the wAddPCA() function
from the Correlplot package (Graffelman and De Leeuw
(\citeproc{ref-graffelman_deleeuw_M_26}{2026})).

\sectionbreak

\section*{References}\label{references}
\addcontentsline{toc}{section}{References}

\protect\phantomsection\label{refs}
\begin{CSLReferences}{1}{1}
\bibitem[\citeproctext]{ref-graffelman_deleeuw_M_26}
Graffelman, Jan, and Jan De Leeuw. 2026. \emph{{Correlplot: A Collection
of Functions for Graphing Correlation Matrices}}.
\url{https://CRAN.R-project.org/package=Correlplot}.

\bibitem[\citeproctext]{ref-groenen_giaquinto_kiers_03}
Groenen, Patrick J. F., Patrizia Giaquinto, and Henk A. L Kiers. 2003.
\emph{{Weighted Majorization Algorithms for Weighted Least Squares
Decomposition Models}}. Econometric Institute Report EI 2003-09.
Econometric Institute, Erasmus University Rotterdam.
\url{https://repub.eur.nl/pub/1700}.

\bibitem[\citeproctext]{ref-gupta_nagar_00}
Gupta, A. K., and D. K. Nagar. 2000. \emph{Matrix Variate
Distributions}. {Chapman \& Hall}.

\bibitem[\citeproctext]{ref-kiers_97}
Kiers, Henk A. L. 1997. {``{Weighted Least Squares Fitting Using
Iterative Ordinary Least Squares Algorithms}.''} \emph{Psychometrika}
62: 251--66.

\bibitem[\citeproctext]{ref-lange_16}
Lange, K. 2016. \emph{MM Optimization Algorithms}. SIAM.

\end{CSLReferences}




\end{document}
